% Cote 119: offprints, reports, courses, CV and letters (typed and printed part): archive cote 119, site page 63 (PDF page 64).
% offprint. Draft compared: none.
% AI-transcribed by Claude Opus 5.5 (claude-opus-5-5, adaptive thinking, effort high), Message Batches API, batch r119_full (msgbatch_01DwZNm5Y7YsKBWyAN4sFs8w), 2026-11-13; not yet checked by eye.
\scanpage{119/63}
\mspage{( 3 )}

2.4. Supposons Pol $(C) \neq \emptyset$. Ceci entraîne que le lien de $C$ est réductif,
et que le $\mathbf{R}$-groupe algébrique affine Z est compact; rappelons qu'un
$\mathbf{R}$-groupe algébrique affine H est dit compact si H $(\mathbf{R})$ est compact et
Zariski-dense dans H. En particulier, on a des isomorphismes canoniques
\[ (2.4.1) \qquad \mathrm{H}^1\,(\mathbf{R}, \mathrm{Z}) \simeq {}_2\mathrm{Z}\,(\mathbf{R}); \]
\[ (2.4.2) \qquad \mathrm{H}^2\,(\mathbf{R}, \mathrm{Z}) \simeq \mathrm{Z}\,(\mathbf{R})/\mathrm{Z}\,(\mathbf{R})^2. \]

Le groupe Z $(\mathbf{R})$ agit sur l'ensemble Pol $(C)$ de la façon suivante : si
\[ \pi \in \mathrm{Pol}\,(C), \quad z \in \mathrm{Z}\,(\mathbf{R}), \quad \mathrm{V} \in \mathrm{ob}\, C, \qquad \varphi \in (z\,\pi)\,(\mathrm{V}) \iff \varphi \circ (\mathrm{id}_{\mathrm{V}} \otimes z_{\mathrm{V}}) \in \pi\,(\mathrm{V}); \]
on a d'autre part
\[ (2.4.3) \qquad \varepsilon\,(z.\pi) = z^2.\varepsilon\,(\pi). \]

D'autre part, le groupe $\mathrm{Aut}_{\mathrm{L}}\,(C)$ des classes d'isomorphisme d'équi-%
valences de $C$ avec elle-même liées par l'identité du lien L de $C$ agit
sur Pol $(C)$ par transport de structure, et on sait par [$(^1)$, 1.6] qu'on a
un isomorphisme canonique $\mathrm{Aut}_{\mathrm{L}}\,(C) \simeq \mathrm{H}^1\,(\mathbf{R}, \mathrm{Z})$. Si on identifie $\mathrm{Aut}_{\mathrm{L}}\,(C)$,
avec ${}_2\mathrm{Z}\,(\mathbf{R})$ par (2.4.1), on vérifie que l'action de ${}_2\acute{\mathrm{Z}}\,(\mathbf{R})$ sur Pol $(C)$
obtenue ainsi coïncide avec l'action déduite de l'action précédente de Z $(\mathbf{R})$.

Exemple 2.5. --- Soient G un $\mathbf{R}$-groupe algébrique affine, $\mathrm{C} \in \mathrm{G}\,(\mathbf{R})$ tel
que $\mathrm{C}^2 \in \mathrm{Z}\,(\mathbf{R})$. On pose $C = Rep_0\,(\mathrm{G})$. Si pour chaque objet V de $C$
il existe une forme de C-polarisation sur V (voir 1.4), l'ensemble de ces
formes définit une polarisation $\pi_{\mathrm{C}}$ de $C$, de parité $\varepsilon\,(\pi_{\mathrm{C}}) = \mathrm{C}^2$ [voir $(^3)$, 5.1].
Un élément C définissant une polarisation du type $\pi_{\mathrm{C}}$ est appelé hodgien,
et une polarisation du type $\pi_{\mathrm{C}}$ est appelée hodgienne. Enfin, G est dit
hodgien s'il existe un élément hodgien dans G $(\mathbf{R})$. On prouve [voir $(^4)$, 2.8]
que C est hodgien si et seulement si la forme $\mathrm{G}_{\mathrm{C}}$ de G définie par la conju-%
gaison $g \mapsto \mathrm{C}\,\bar{g}\,\mathrm{C}^{-1}$ sur G $(\mathbf{C})$ est compacte (au sens de 2.4). Si C, C$'$ sont
des éléments hodgiens, on a $\pi_{\mathrm{C}} = \pi_{\mathrm{C}'}$ si et seulement s'il existe $z \in \mathrm{Z}\,(\mathbf{R})$
(nécessairement unique), $g \in \mathrm{G}\,(\mathbf{R})$ tels que $\mathrm{C}' = zg\,\mathrm{C}\,g^{-1}$. Enfin, G est
hodgien si et seulement s'il est réductif et est une forme tordue intérieure
de sa forme compacte K; ceci revient aussi à dire que le lien de G est le
lien d'un $\mathbf{R}$-groupe compact.

3. Classification des polarisations; liens polarisables.

Proposition 3.1. --- L'ensemble Pol $(C)$ est un pseudo-torseur (i. e. vide
ou un ensemble principal homogène) sous l'action de Z $(\mathbf{R})$ (2.4); de même,
si $\varepsilon \in \mathrm{Z}\,(\mathbf{R})$, $\mathrm{Pol}_{\varepsilon}\,(C)$ est un pseudo-torseur sous ${}_2\mathrm{Z}\,(\mathbf{R})$.

3.2. Il résulte de cette proposition et de (2.4.3) que pour $\pi \in \mathrm{Pol}\,(C)$
l'image de $\varepsilon\,(\pi)$ dans $\mathrm{Z}\,(\mathbf{R})/\mathrm{Z}\,(\mathbf{R})^2$ est indépendante de $\pi \in \mathrm{Pol}\,(C)$ et déter-%
mine donc, par (2.4.2), une classe de cohomologie $\varepsilon_C \in \mathrm{H}^2\,(\mathbf{R}, \mathrm{Z})$ chaque fois
que Pol $(C) \neq \emptyset$. On peut interpréter $\varepsilon_C$ comme l'obstruction à l'existence
sur $C$ d'une polarisation symétrique (2.2).

\mspage{\hand{63}}
